\documentclass[12pt,a4paper]{jsarticle}

% ==================== パッケージ ====================
\usepackage[utf8]{inputenc}
\usepackage{amsmath,amssymb}
\usepackage{graphicx}
\usepackage{booktabs}
\usepackage{multirow}
\usepackage{url}
\usepackage{cite}
\usepackage{algorithm}
\usepackage{algpseudocode}
\usepackage[dvipdfmx]{hyperref}
\usepackage{setspace}
\usepackage{geometry}

\geometry{
  top=25mm, bottom=25mm,
  left=25mm, right=25mm
}

\setstretch{1.3}

% ==================== タイトル ====================
\title{
  \Large 所有者優先権を保証するプリエンプティブ方式を用いた\\
  学内GPU共有基盤の提案とユーザ参加行動の評価
  \vspace{1em}
}
\author{
  櫻井 翔琉\\
  大阪大学 工学部\\
  指導教員：大下 裕一，下西 英之
}
\date{2026年3月}

% ==================== 本文 ====================
\begin{document}

\maketitle
\tableofcontents
\newpage

% ============================================================
\section{はじめに}
% ============================================================

近年，大規模言語モデル（LLM）や生成AIの急速な発展は，GPUを中心とする計算資源への需要を爆発的に増大させている\cite{lecun2015deep,brown2020language}．
ChatGPTやStable Diffusionに代表されるAIシステムの訓練・推論には膨大な並列演算処理が必要であり，高性能GPUはもはや現代の研究・開発における不可欠な基盤となっている．
しかしながら，NVIDIA A100 等の最新GPUは1基あたり数十万円から数百万円に及ぶ高額な機器であり，供給不足の慢性化も相まって，個々の研究室や部門が最新の計算環境を独自に整備・維持することは経済的に困難な状況にある\cite{hooker2021hardware}．

この状況は大学等の研究機関において特に顕著である．
組織内には高性能から低性能まで多様なGPUが各研究室に分散して存在するものの，その多くは所有研究室の作業スケジュールに依存した断続的な使用に留まっており，
深夜・週末・実験間インターバルといった時間帯には大量のアイドル時間が発生していることが報告されている\cite{jeon2019analysis}．
一方で，低性能なGPUしか保有していない研究室では，タスクの完了まで数十万秒を超える待ち時間が発生し，実質的に計算不能な状態（飽和状態）に陥る場合も少なくない．
このような組織内の資源格差は，研究活動の進展を大きく阻害する要因となっている\cite{armbrust2010view,buyya2009cloud}．

こうした課題を解決するアプローチとして，学内に分散するGPU資源をネットワークを介して仮想的に統合し，ユーザ間で相互に利用可能とする「学内GPU共有基盤」の構想がある．
本研究ではこれを \textbf{Air Computing Pool（ACP）}と称する．
ACPの基本的なアイデアは，個々の資源における「遊休時間」に着目し，その時間帯に他者のタスクを動的に割り当てることで，組織全体の計算スループットを最大化することにある．
これはGrid Computingにおいて試みられたコミュニティ内資源共有と類似した概念であるが，ACPはボランティア精神に依存するのではなく，
参加者が\textbf{自発的に参加し続けるインセンティブメカニズム}を内包する点において本質的に異なる．

ACPの設計において最も重要な課題は，\textbf{資源提供者（GPU所有者）の参加インセンティブを確保する}ことである．
一般的なスケジューリング手法，例えばFCFS（先着順）や支配的資源公平分配（DRF）は，すべてのユーザを等しく扱うことを前提としている\cite{ghodsi2011dominant}．
これをACPに適用した場合，高性能GPUを提供した所有者が自身のタスクを実行しようとした際に，
他者のゲストタスクが当該GPUを占有中であれば，所有者であってもその完了を待機しなければならない．
この状況では，資源共有によって他者の高性能資源を利用できる利益よりも，自身が提供した資源を即座に利用できないという不利益の方が大きくなる場合があり，
高性能GPU所有者のACP離脱を招く．
高性能資源が失われると，それに依存していた低性能ユーザの恩恵も消滅し，最終的にシステム全体が崩壊するという\textbf{負の連鎖}が生じる\cite{feldman2004robust}．

本研究では，この根本的な課題を解決するために，\textbf{所有者優先権を保証するプリエンプティブ方式}を採用した学内GPU共有基盤を提案する．
本方式では，所有者が自身のGPUにタスクを投入した際，実行中のゲストタスクを即座に強制中断（プリエンプション）させ，所有者に資源を即時返却する．
これにより，所有者は「自身の資源はいつでも利用可能」という技術的保証を得た上で，他者の遊休資源も利用可能となる．
結果として，ACPへの参加はいかなる場合にも参加者にとって有利となり，システムが自律的に安定状態へと収束することが期待される．

本研究の主要な貢献は以下の3点である．

\begin{enumerate}
  \item \textbf{技術的保証の導入：}
    所有者がGPUを必要とした際，実行中のゲストタスクを即座に中断・退避させる制御機構を導入し，
    所有者の待ち時間を物理的に排除する技術的保証（所有者優先権）を与えた．

  \item \textbf{TAT短縮の実証：}
    プリエンプションの導入により，提供者は自資源の即時利用権を維持しながら他者の遊休資源も並行活用可能となり，
    単独利用時よりタスク完了時間（TAT）が短縮されることをシミュレーションにより定量的に実証した．
    特に高性能GPU所有者のTATを非共有時と同等以下に抑えることに成功した．

  \item \textbf{持続可能性の実証：}
    複雑な経済的報酬や金銭的インセンティブを必要とせず，プリエンプションという技術的保証のみで
    システムの持続可能性（全ユーザの参加継続率100\%）を確保できることを参加行動モデルを用いて示した．
\end{enumerate}

本論文の構成は以下の通りである．
第2章では関連研究との位置づけを整理する．
第3章では提案するACPのシステムモデルと数理モデルを詳述する．
第4章ではシミュレーションによる性能評価と持続性評価の結果を示す．
第5章では評価結果の考察と議論を行い，第6章で本研究の結論と今後の課題を述べる．

% ============================================================
\section{関連研究}
% ============================================================

\subsection{GPU共有基盤に関する既存研究}

大規模な計算資源の共有・管理に関する研究は，クラスタコンピューティングの黎明期から継続的に行われてきた．
データセンター向けのクラスタ管理システムとしては，KubernetesやSLURMが広く実用されており，
ジョブキューイングとリソース割り当てを統合的に管理する\cite{armbrust2010view}．
しかしこれらのシステムは，資源を組織が一元的に所有・運用することを前提としており，
\textbf{個人が所有する分散資源を共有する}という本研究の設定とは本質的に異なる．

商業クラウドサービス（AWS，GCP，Azure等）は，GPUを時間単位でオンデマンドに提供するが，
コストが高く，学内に既存する資源の活用という観点からは補完的な役割に限定される\cite{buyya2009cloud}．
また，GPU購入を需要者と所有者がマッチングするプラットフォームも存在するが，
金銭的報酬を介した市場メカニズムに依存しており，
学内コミュニティの非商業的な相互扶助という文脈には馴染まない．

\subsection{スケジューリングと公平分配}

公平資源配分の分野では，Dominant Resource Fairness (DRF)\cite{ghodsi2011dominant}が代表的な研究である．
DRFは複数種類の資源を考慮した公平分配を実現するが，そのアプローチはすべての参加者を等しく扱うことを前提としており，
資源提供者に対する特別な優先権を付与する仕組みを持たない．
本研究のように，\textbf{提供量に応じた優先権}を保証するというアプローチとは目的が異なる．

P2Pネットワークにおける参加インセンティブ設計の研究も本研究と関連する．
Feldmanら\cite{feldman2004robust}は，P2Pシステムにおける「フリーライダー問題」を分析し，
貢献に基づいた報酬メカニズムの設計が参加維持に不可欠であることを示した．
本研究の問題意識はこれと共通するが，本研究では金銭的・評判的報酬ではなく，
\textbf{技術的保証（プリエンプション）}によってインセンティブを担保する点が根本的に異なる．

\subsection{GPUプリエンプション技術}

GPU上でのプリエンプションに関する研究は，主にシステムのスループット最大化を目的として行われてきた．
複数ジョブのGPU上での共存実行（co-location）や，メモリの動的割り当てによる利用効率向上，
仮想化環境においてGPU資源を複数の仮想マシンに分割して提供する機構など，
様々な技術的アプローチが提案されている\cite{jeon2019analysis}．

これらの既存研究に共通するのは，プリエンプションを\textbf{スループット最大化のための効率化ツール}として位置づけている点である．
これに対して本研究では，プリエンプションを\textbf{所有者の参加インセンティブを維持するための所有権保証メカニズム}として再解釈する．
すなわち，プリエンプションの技術的機能（タスクの即座中断・再開）を活用しながら，
その目的を「効率化」から「所有権の物理的保証」へと転換することが本研究の新規性の核心である．

\subsection{本研究の位置づけ}

表\ref{tab:related}に，関連研究との比較を示す．
本研究は，GPU共有基盤としてのシステム設計と，参加インセンティブ維持のためのメカニズム設計を統合して扱う点，
そして複雑な経済的報酬機構なしにシステムの持続可能性を保証する点において，
既存研究に対する独自の貢献を持つ．

\begin{table}[h]
  \centering
  \caption{関連研究との比較}
  \label{tab:related}
  \begin{tabular}{lcccc}
    \toprule
    アプローチ & 所有権保証 & 共有基盤 & インセンティブ & 非金銭的 \\
    \midrule
    FCFS / DRF & $\times$ & $\circ$ & $\times$ & --- \\
    P2P報酬機構 & $\times$ & $\circ$ & $\circ$ & $\times$ \\
    クラウドサービス & $\circ$ & $\circ$ & $\times$ & $\times$ \\
    既存GPU Preemption & $\times$ & $\circ$ & $\times$ & --- \\
    \textbf{本研究（ACP）} & $\mathbf{\circ}$ & $\mathbf{\circ}$ & $\mathbf{\circ}$ & $\mathbf{\circ}$ \\
    \bottomrule
  \end{tabular}
\end{table}

% ============================================================
\section{システムモデル}
% ============================================================

\subsection{学内GPU共有基盤ACPの概要}

我々が提案する学内GPU共有基盤「Air Computing Pool（ACP）」は，大学の研究室等に点在する計算資源をネットワークを介して統合し，ユーザ間での相互利用を可能にする基盤である．
共有対象は，各研究室が個別に導入・運用しているGPUサーバであり，所有者は当該資源を最優先に使用する権利を保持する．
ACPは，これら個別資源の遊休時間を他者のタスクへ動的に割り当てることで，組織全体の計算スループット最大化を図ることを目的とする．

運用において，所有者は提供可能な時間帯を任意に設定でき，利用者はACPの管理ポータル上に表示される完了予測時間に基づいて最適なサーバを選択する．
タスクの割り当てはシステムが自動で行い，ユーザは複雑な設定なしにクラスタ内の最適な計算資源を利用できる．

\subsection{所有者優先権を保証するプリエンプティブ方式}

研究活動における計算ニーズは突発的であるため，所有者が遊休時間を事前に正確に把握することは困難である．
タスク終了まで資源が返却されない非プリエンプティブな方式では，所有者の即時利用が損なわれる懸念があり，これが共有への参加障壁となる．

そこで本研究では，所有者の即時利用権を技術的に保証するため，ゲストタスクを強制排除するプリエンプティブ方式を採用する．
具体的には，ノード内でのゲストジョブ実行中に所有者のジョブ投入を検知すると，システムはゲストジョブを強制的に中断・退避させる．
この際，進捗損失を防ぐため，未処理タスクをチェックポイントとして保存し，マイグレーションを行う．
中断されたゲストジョブは，以下の3戦略から最短の完了が見込まれる先を選択して再開される．

\begin{enumerate}
  \item 自身の所有するGPUへの復帰
  \item 元のGPUにおける所有者タスク終了までの待機
  \item 他の空きGPUへの再投入
\end{enumerate}

\subsection{モデル化と記号定義}

提案方式がシステム全体および各ユーザに与える影響を定量的に評価するため，待ち行列理論に基づいた数理モデルを構築する．
各構成要素の定義および数理モデルで使用する記号を表\ref{tab:notation}に示す．

\begin{table}[h]
  \centering
  \caption{モデルの構成要素と記号定義}
  \label{tab:notation}
  \begin{tabular}{ll}
    \toprule
    記号 & 定義・意味 \\
    \midrule
    $U$ & システムを利用する全ユーザの集合 \\
    $u \in U$ & 個別のユーザ（自身のGPU $g_u \in G$ を所有） \\
    $\lambda_u$ & ユーザ $u$ によるタスクの平均到着率 \\
    $S_j$ & ジョブ $j$ の仕事量（タスクサイズ，単位：TFLOP） \\
    $G$ & システム内の全GPUの集合 \\
    $\mu_g$ & GPU $g$ が持つ固有の処理性能（単位：TFLOP/s） \\
    $T_\mathrm{pred}$ & ジョブ投入時に算出される完了予測時間 \\
    $TAT$ & ターンアラウンドタイム（到着〜完了の経過時間） \\
    $\rho_\mathrm{own}$ & 所有者の資源占有率（稼働率） \\
    $\mu_\mathrm{eff}$ & ゲストから見た実効処理性能 \\
    $\alpha$ & プリエンプション・オーバーヘッド係数 \\
    $S_\mathrm{rem}$ & 中断されたタスクの残仕事量 \\
    \bottomrule
  \end{tabular}
\end{table}

\subsection{比較対象のスケジューリング方式}

本研究では，以下の4つのスケジューリング方式を比較対象とする．

\paragraph{シナリオ1：非共有（No Sharing）}
従来の研究機関の現状であるACPを導入しないシナリオである．
各ユーザが自身の所有するGPUのみを使用する．
完了予測時刻は，自身のGPUに現在たまっているタスクの完了予測時刻に，自身のタスク実行時間を加えたものとなる．

\begin{equation}
  T_\mathrm{pred} = T_\mathrm{finish} + \sum_{j \in Q_\mathrm{self}} \frac{S_j}{\mu_\mathrm{self}} + \frac{S_\mathrm{new}}{\mu_\mathrm{self}}
  \label{eq:no_sharing}
\end{equation}

ここで，$T_\mathrm{finish}$は実行中タスクの終了時刻，$Q_\mathrm{self}$はキュー内のタスク集合，$\mu_\mathrm{self}$は自機の処理性能を表す．

\paragraph{シナリオ2：FCFS（先着順）}
全ユーザが全GPUを利用可能とするが，優先度制御やプリエンプションは行わないシナリオである．
ジョブ投入時に全ノードの$T_\mathrm{pred}$を計算し，最小となるノードを選択する．

\begin{equation}
  T_\mathrm{pred} = \min_{g \in G} \left( T_{\mathrm{finish},g} + \sum_{j \in Q_g} \frac{S_j}{\mu_g} + \frac{S_\mathrm{new}}{\mu_g} \right)
  \label{eq:fcfs}
\end{equation}

\paragraph{シナリオ3：所有者優先方式（Owner Priority）}
FCFS に対し，リソース所有者の優先権を部分的に導入したシナリオである．
実行中のタスクを中断するプリエンプションは行わないが，所有者のタスクが投入された際には，
待機キュー内の他者タスクを追い越して先頭へ挿入される優先制御を行う．

他者のGPU（ゲスト利用）を選択する場合，将来的に発生する所有者タスクの割り込みによる「実効的な処理性能の低下」を考慮する必要がある．
所有者が資源を占有する割合を示す稼働率$\rho_\mathrm{own}$を次式で定義する．

\begin{equation}
  \rho_\mathrm{own} = \frac{\lambda_\mathrm{own} \cdot \bar{S}_\mathrm{own}}{\mu_g}
  \label{eq:rho}
\end{equation}

ここで$\lambda_\mathrm{own}$は所有者のタスク発生率，$\bar{S}_\mathrm{own}$は平均タスクサイズ，$\mu_g$はGPUの基本処理レートである．
この$\rho_\mathrm{own}$を用いて，ゲストから見た実効処理性能$\mu_\mathrm{eff}$を次のように定式化する．

\begin{equation}
  \mu_\mathrm{eff} = \frac{\mu_g}{1 + \rho_\mathrm{own}}
  \label{eq:mu_eff}
\end{equation}

ユーザ$u$は，自機利用時（所有権あり）と他機利用時（ゲスト）で以下の式を使い分け，最短のGPUを選択する．

自機選択時（他者タスクを無視し，自身の先行タスクのみを考慮）：
\begin{equation}
  T_{\mathrm{pred,own}} = T_{\mathrm{finish},g} + \sum_{j \in Q_{g,\mathrm{self}}} \frac{S_j}{\mu_g} + \frac{S_\mathrm{new}}{\mu_g}
  \label{eq:own_select}
\end{equation}

他機選択時（実効性能$\mu_\mathrm{eff}$を用いて全タスクの処理時間を算出）：
\begin{equation}
  T_{\mathrm{pred,other}} = T_{\mathrm{finish},g} + \sum_{j \in Q_g} \frac{S_j}{\mu_\mathrm{eff}} + \frac{S_\mathrm{new}}{\mu_\mathrm{eff}}
  \label{eq:other_select}
\end{equation}

\paragraph{シナリオ4：プリエンプティブ方式（Preemptive）}
本研究の提案手法であり，所有者の即時資源利用を保証する．
所有者投入時，実行中のゲストタスクは即座に中断され，待ち時間は自身の前タスクが完了するまでとなる．
ゲストタスクの完了予測時間には途中切断リスクが生じるため，予測式に期待ペナルティ$E[\mathrm{Penalty}]$を算入する．

\begin{equation}
  T_{\mathrm{pred,other}} = T_{\mathrm{finish},g} + \sum_{j \in Q_g} \frac{S_j}{\mu_\mathrm{eff}} + \frac{S_\mathrm{new}}{\mu_\mathrm{eff}} + E[\mathrm{Penalty}]
  \label{eq:preemptive_pred}
\end{equation}

\begin{equation}
  E[\mathrm{Penalty}] = (\lambda_\mathrm{own} \cdot T_\mathrm{svc}) \cdot (\bar{T}_{\mathrm{svc,own}} \cdot (1 + \alpha))
  \label{eq:penalty}
\end{equation}

ここで$\alpha$は中断・再開に伴うオーバーヘッド係数（本研究では$\alpha = 0.2$）である．

中断されたゲストタスクは，残仕事量$S_\mathrm{rem}$を保持し，以下の3つの選択肢から予測完了時刻が最小となるものを選択して再開する．

\begin{equation}
  \min \begin{cases}
    T_\mathrm{fin} + \sum \frac{S_j}{\mu_\mathrm{self}} + \frac{S_\mathrm{rem}}{\mu_\mathrm{self}} & \text{（自機復帰）} \\
    T_\mathrm{fin} + \frac{S_\mathrm{own}}{\mu_\mathrm{own}} + \frac{S_\mathrm{rem}}{\mu_\mathrm{eff}} + E[\mathrm{Pnlty}] & \text{（元機待機）} \\
    \min_{g' \in G} \left( T_{\mathrm{pred,other}}(S_\mathrm{rem}) \right) & \text{（他機移行）}
  \end{cases}
  \label{eq:resume}
\end{equation}

\subsection{保護比率（Protection Ratio）の定義}

各スケジューリング方式が所有者の利益をどの程度保護できているかを定量的に評価するために，\textbf{保護比率（Protection Ratio; PR）}を次式で定義する．

\begin{equation}
  \mathrm{PR}_i = \frac{TAT_{\mathrm{shared},i}}{TAT_{\mathrm{NoSharing},i}}
  \label{eq:pr}
\end{equation}

ここで$TAT_{\mathrm{shared},i}$は共有環境下でのユーザ$i$の平均TATであり，$TAT_{\mathrm{NoSharing},i}$は非共有時の平均TATである．
$\mathrm{PR} \leq 1$であれば「共有することで損をしない（参加が合理的）」ことを意味し，
$\mathrm{PR} > 1$であれば「共有することで損をする（離脱の動機が生じる）」ことを意味する．

\subsection{参加意思決定モデルの定式化}

実環境では，自タスクの遅延が許容限度を超えた場合，ユーザがリソース提供を停止するリスクがある．
そこで，ユーザ$i$を自身の平均TATに基づいて行動する合理的エージェントと定義し，以下のモデルを構築した．

各フェーズにおけるユーザ$i$の効用$U_i$を，単独利用時の完了時間$TAT_{\mathrm{std},i}$と共有環境下での観測完了時間$TAT_{\mathrm{obs},i}$の比として定義する．

\begin{equation}
  U_i = \frac{TAT_{\mathrm{std},i}}{TAT_{\mathrm{obs},i}}
  \label{eq:utility}
\end{equation}

ここで，$TAT_{\mathrm{obs},i}$は，参加ユーザについては前フェーズの実測値を，離脱ユーザについてはパイロットテストによる予測値を用いる．
ユーザ$i$は，次フェーズにおける意思決定$D_i \in \{\text{参加}, \text{離脱}\}$を以下の判定式に従って決定する．

\begin{equation}
  D_i = \begin{cases}
    \text{参加} & (U_i \geq 1) \\
    \text{離脱} & (U_i < 1)
  \end{cases}
  \label{eq:decision}
\end{equation}

保護比率の定義（式\ref{eq:pr}）と照合すると，$\mathrm{PR}_i \leq 1$のとき$U_i \geq 1$となり参加が継続されることがわかる．
すなわち，\textbf{プリエンプティブ方式が全負荷においてPR $\leq$ 1を保証できれば，全ユーザの参加継続が自動的に達成される}．

% ============================================================
\section{評価}
% ============================================================

\subsection{シミュレーション環境}

\subsubsection{ユーザおよびGPU構成}

ユーザ数は18名とし，表\ref{tab:users}に示す通り9段階の計算性能ティアに各ティア2名ずつ割り当てる．
ティアの区分は，実際の大学研究室環境で観測される多様なGPU性能分布を反映するよう設計した．
ティア1（GTX 1650）からティア9（A100）まで，最大約100倍以上の性能差を持つ異種資源が混在する環境を想定する．

\begin{table}[h]
  \centering
  \caption{シミュレーションにおけるユーザ構成とGPU性能パラメータ}
  \label{tab:users}
  \begin{tabular}{llllr}
    \toprule
    ティア & ユーザー数 & ユーザーID & 想定GPUモデル & 計算性能 $\mu$ \\
    \midrule
    Tier 1 & 2 & 0, 9  & GTX 1650     &   2.98 \\
    Tier 2 & 2 & 1, 10 & GTX 1080     &   8.87 \\
    Tier 3 & 2 & 2, 11 & RTX 3050     &  20.41 \\
    Tier 4 & 2 & 3, 12 & RTX 3060 Ti  &  64.83 \\
    Tier 5 & 2 & 4, 13 & RTX 2080     &  82.60 \\
    Tier 6 & 2 & 5, 14 & RTX 2080 Ti  & 110.00 \\
    Tier 7 & 2 & 6, 15 & Titan RTX    & 180.50 \\
    Tier 8 & 2 & 7, 16 & RTX 4070     & 233.00 \\
    Tier 9 & 2 & 8, 17 & A100         & 311.84 \\
    \bottomrule
  \end{tabular}
\end{table}

\subsubsection{タスク生成モデル}

タスクの到着過程はポアソン過程に従うものとし，到着率$\lambda = 0.005$\,[tasks/s]（全ユーザ均一）とする．
タスクサイズは，推論タスクと学習タスクの2種類を混合して生成する．
推論タスクは比較的小さいサイズ（平均：9{,}580 TFLOP），学習タスクは大きいサイズ（平均：412{,}180 TFLOP）とし，
学習タスクの割合（training ratio）を評価条件として変化させる．

基本設定（Uniform Workload）では全ユーザの training ratio を0.3とし，
計算量は1画像あたりの基準計算量$S_\mathrm{base} = 0.360$\,[TFLOP]，エポック数$E = 10$，
バッチサイズをグループAは3{,}000枚，グループBは6{,}000枚として生成する．
総シミュレーション時間は86{,}400\,s（1日分）とする．

\subsubsection{システム負荷率の定義}

システム全体の負荷率$\rho_\mathrm{sys}$を次式で定義する．

\begin{equation}
  \rho_\mathrm{sys} = \frac{\text{総発生タスク量}}{\text{全GPUの総処理能力}}
  \label{eq:sys_load}
\end{equation}

負荷率は0.1から1.0まで0.1刻みで変化させ，各条件において10試行の平均を評価値とする．
バッチサイズを調整することで$\rho_\mathrm{sys}$を制御する．

\subsection{均一ワークロード評価}

\subsubsection{全体平均TATの比較}

図\ref{fig:overall_tat}に，全ユーザの平均TATをシステム負荷率ごとに示す（対数スケール）．

非共有方式では，負荷率0.3以上から平均TATが急激に増大し，高負荷域では他の共有方式と比較して100〜1000倍以上の値を示す．
これは，低性能GPU所有者（Tier 1〜3）が自身のGPUだけでは処理しきれず，飽和状態に陥るためである．

共有3方式（FCFS，所有者優先，プリエンプティブ）の比較では，
低〜中負荷域（$\rho_\mathrm{sys} \leq 0.7$）においてプリエンプティブ方式が全シナリオ中で最も低いTATを維持している．
ただし，$\rho_\mathrm{sys} = 0.9$以上の高負荷域ではプリエンプティブ方式のTATが上昇し，
所有者優先方式を上回る場合が確認される．
この逆転現象の原因は，低・中性能GPU群のゲスト利用時に頻発するプリエンプションによって，
タスクの中断・再投入サイクルが繰り返され，ペナルティが蓄積するためである．

\subsubsection{ティア別TATの詳細分析}

各ユーザIDの平均TATをシナリオ別に示したものが表\ref{tab:tat_detail}である．

\begin{table}[h]
  \centering
  \caption{各ユーザのシナリオ別平均TAT詳細}
  \label{tab:tat_detail}
  \small
  \begin{tabular}{lrrrr}
    \toprule
    ユーザID & 非共有 & FCFS & 所有者優先 & プリエンプティブ \\
    \midrule
    0  & 830{,}915 s & 314 s & 206 s & 230 s \\
    1  & 224{,}255 s & 350 s & 191 s & 241 s \\
    2  &  66{,}955 s & 321 s & 204 s & 238 s \\
    3  &     588 s   & 380 s & 228 s & 190 s \\
    4  &     299 s   & 497 s & 269 s & 171 s \\
    5  &     282 s   & 338 s & 198 s & 146 s \\
    6  &      78 s   & 356 s & 180 s &  79 s \\
    7  &      58 s   & 435 s & 147 s &  58 s \\
    8  &      40 s   & 383 s & 156 s &  40 s \\
    9  & 1{,}450{,}533 s & 483 s & 379 s & 342 s \\
    10 &   490{,}566 s & 501 s & 371 s & 360 s \\
    11 &   201{,}421 s & 578 s & 466 s & 358 s \\
    12 &    30{,}591 s & 533 s & 311 s & 324 s \\
    13 &    11{,}810 s & 460 s & 318 s & 303 s \\
    14 &     4{,}053 s & 461 s & 289 s & 266 s \\
    15 &       243 s   & 644 s & 300 s & 191 s \\
    16 &       166 s   & 723 s & 236 s & 168 s \\
    17 &       117 s   & 418 s & 231 s & 114 s \\
    \bottomrule
  \end{tabular}
\end{table}

Tier 1〜3に属する低性能GPU所有者（ユーザ0, 1, 2, 9, 10, 11）の非共有シナリオでは，
仕事量に対してGPU性能が著しく不足しているため，TATが数十万秒から数百万秒という非常に高い値を示し，実質的な計算不能状態（飽和状態）となっている．
これに対し，いずれの共有シナリオにおいても，これらのユーザのTATは数百秒程度へと大幅に改善されている．
ACPを導入することで低性能な自機資源の不足を補い，クラスタ内の高性能な計算資源を活用可能になったためである．

Tier 8〜9（ユーザ7, 8, 16, 17）の高性能GPU所有者に着目すると，スケジューリングポリシーによる差が顕著に現れる．
FCFS・所有者優先シナリオでは，これらのユーザのTATは非共有時と比較して数倍から十数倍へと増加している．
自機の高性能GPUが他者のタスクによって占有され，自身のタスクが長時間待機を強いられているためである．
対照的に，提案手法であるプリエンプティブ方式では，高性能GPU所有者のTATは非共有時と同等水準に抑えられており，
所有者のジョブ投入時にゲストタスクを即座に排除する所有者優先権の保証が有効に機能していることが確認できる．

\subsection{保護比率分析}

\subsubsection{高性能ユーザにおけるPRの推移}

図\ref{fig:protection_ratio}に，Tier 9ユーザの保護比率（PR）をシステム負荷率ごとに示す．

FCFS方式では，$\rho_\mathrm{sys} \geq 0.8$からPR $> 1$となり，高負荷時（$\rho_\mathrm{sys} = 1.0$）ではPRが4.96に達する．
これは，A100所有者が「ACP に参加することで非共有時より約5倍もTATが悪化する」状態を意味し，
合理的なユーザであれば参加を継続する動機が失われることを示している．

所有者優先方式では，FCFS よりも悪化は緩和されるが，$\rho_\mathrm{sys} \geq 0.9$においてPR $> 1$となる．
実行中タスクを停止できないためにゲストタスク完了まで待機が生じ，高負荷時に所有者保護が不完全となる．

プリエンプティブ方式では，全負荷帯（$\rho_\mathrm{sys} = 0.1 \sim 1.0$）を通じてPR $\leq 1$を維持する．
これは，所有者がタスクを投入した瞬間にリソースを即時奪還できるプリエンプションの特性が，
周囲の混雑状況に関わらず「所有者の優先権」を物理的に保証しているためである．

\subsubsection{PRの数理的解釈}

式(\ref{eq:pr})で定義したPRと，式(\ref{eq:decision})の意思決定モデルとの対応関係から，
PR $\leq 1$はすなわち「ACP参加によってTATが悪化しない」ことを意味する．
プリエンプティブ方式が全負荷でPR $\leq 1$を保証できる理由は，
所有者がゲストタスクを即座に排除することで，自身の待ち時間が非共有時のキュー待ち以下に抑えられるためである．

\subsection{持続性評価}

\subsubsection{反復型シミュレーションの設定}

前節までの評価は全ユーザの常時参加を前提としていたが，
実環境ではTATの悪化が許容限度を超えた場合，ユーザがリソース提供を停止するリスクがある．
本節では，式(\ref{eq:utility})・(\ref{eq:decision})で定義した参加意思決定モデルを用いた反復型シミュレーションにより，
各スケジューリング方式がシステムの持続性に与える影響を評価する．

シミュレーション環境は第4.1節の18名・9ティア構成を維持し，
初期状態では全ユーザに対して50\%の確率でランダムにACPへ参加させた状態から開始する．
各イテレーションで参加状況と観測TATを更新し，10イテレーション分の推移を観察する．

\subsubsection{参加者数の推移と各方式の挙動}

図\ref{fig:sustainability}に，シナリオごとの参加者数の推移を示す．

\paragraph{FCFS方式}
所有者保護のない本方式では，高性能層が第1イテレーション直後のTAT悪化により，
第3イテレーションまでに全員離脱した．
有力な計算資源の消失は中・低性能層の効用低下を招き，最終的にシステム全体が破綻する\textbf{負の連鎖}が確認された．

\paragraph{所有者優先方式}
優先権の付与により一定の引き止め効果は見られたが，実行中タスクを中断できない待機リスクから
高性能層が第3イテレーションで離脱した．中性能層も不安定な挙動を示し，限定的な持続性に留まった．

\paragraph{プリエンプティブ方式（提案方式）}
本方式は所有者の即時資源利用を物理的に保証するため，高性能層を含む全ユーザの離脱を完全に抑制した．
第2イテレーション以降は参加率100\%の安定状態を維持しており，
ACPの持続においてプリエンプションが不可欠な要素であることが実証された．

\subsubsection{システム負荷率と最終参加者数の関係}

次に，システム全体の負荷状況がユーザの意思決定に与える影響を評価するため，
最終イテレーションにおける参加者数を各負荷率で比較する（図\ref{fig:final_participants}）．

プリエンプティブ方式はほぼすべての負荷率において全18名の参加を維持しており，他シナリオと比較して高い持続性を示している．
ただし，負荷率0.8以降では高性能ユーザ2名が離脱するケースが確認される．
これは，自資源が飽和することで他資源を利用した際の途中切断オーバーヘッドにより，
単独利用時よりわずかにTATが悪化するためである．
それでも，どの負荷率においてもプリエンプティブ方式が他方式を大きく上回る持続性を示している．

\subsection{不均一ワークロード評価}

\subsubsection{評価目的と実験設計}

前節まではすべてのユーザのtraining ratioを0.3に固定した均一ワークロードを評価した．
しかし実環境では，研究内容によってユーザごとのワークロード特性は大きく異なる．
本節では，ティアごとに異なるworkload比率を設定した不均一ワークロード評価を行い，
提案方式の頑健性を検証する．

表\ref{tab:hetero_scenarios}に評価するシナリオを示す．

\begin{table}[h]
  \centering
  \caption{不均一ワークロード評価シナリオ}
  \label{tab:hetero_scenarios}
  \begin{tabular}{lcccp{4.5cm}}
    \toprule
    シナリオ & Low Tier & Mid Tier & High Tier & 評価のねらい \\
    \midrule
    Uniform   & 0.3 & 0.3 & 0.3 & ベースライン \\
    Low-Heavy & 0.7 & 0.3 & 0.1 & 低性能層が大きい学習タスクを多数発生させる \\
    High-Heavy & 0.1 & 0.3 & 0.7 & 高性能層が大きい学習タスクを多数発生させる（論文の核心シナリオ） \\
    Random    & ランダム & ランダム & ランダム & 現実の多様な環境の再現 \\
    \bottomrule
  \end{tabular}
\end{table}

各シナリオは10試行（seed固定）×負荷率0.1〜1.0の条件で実行し，平均・最小・最大値を評価値とする．

\subsubsection{Uniform（ベースライン）}

第4.2・4.3節の結果と一致する．低性能層はACPで劇的にTATが改善し（非共有時と比較して最大約300倍短縮），
高性能層はFCFSで$\rho_\mathrm{sys} = 0.8$以上においてPR $> 1$となる一方，
プリエンプティブ方式は全負荷でPR $\leq 1$を維持する．

\subsubsection{Low-Heavy（低性能層が学習タスク重視）}

低性能層のユーザが大規模な学習タスクを多数発生させるシナリオである．
このシナリオでは低性能層が高性能GPUにゲストとしてアクセスする頻度が高まり，
ACPによる恩恵が均一ワークロードより更に大きくなる．
高性能ユーザへの影響は，低性能ユーザのタスクがプリエンプションによって排除されるため，
均一ワークロードと比較して軽微に留まる．
プリエンプティブ方式は全負荷でPR $\leq 1$を維持することが確認される．

\subsubsection{High-Heavy（高性能層が学習タスク重視）}

本研究の核心シナリオである．A100等の高性能GPU所有者が大規模な学習タスクを頻繁に実行する状況下では，
FCFSにおけるPRの悪化が最も顕著となる．
所有者が長い学習タスクを実行中にゲストが当該GPUを占有し，所有者が長時間待機を強いられるためである．
このシナリオにおいてFCFSのPRは高負荷域で最大値をとる一方，
プリエンプティブ方式では所有者タスク投入時に即座にゲストを排除するため，PR $\leq 1$を維持する．
この結果は，プリエンプティブ方式の優位性がワークロードの不均一性に対して頑健であることを示している．

\subsubsection{Random（ランダム比率）}

各試行においてすべてのユーザのtraining ratioをランダムに再サンプルすることで，
現実の多様なワークロード環境を模擬したシナリオである．
10試行の平均・帯域でのPR評価の結果，プリエンプティブ方式は全シナリオ中で唯一，全負荷帯においてPR $\leq 1$を安定して維持する．

\subsubsection{不均一ワークロード評価のまとめ}

表\ref{tab:hetero_summary}に，4シナリオをまたいだPRの比較をまとめる．

\begin{table}[h]
  \centering
  \caption{各ワークロードシナリオにおけるTier 9 保護比率（PR）の概要}
  \label{tab:hetero_summary}
  \begin{tabular}{lcccc}
    \toprule
    シナリオ & FCFS の最大PR & Owner Priority の最大PR & Preemptive の最大PR & 判定 \\
    \midrule
    Uniform     & $> 1$（load $\geq 0.8$） & $> 1$（load $\geq 0.9$） & $\leq 1$（全負荷） & Preemptive \checkmark \\
    Low-Heavy   & $> 1$                   & $> 1$                    & $\leq 1$           & Preemptive \checkmark \\
    High-Heavy  & 最大悪化               & $> 1$                    & $\leq 1$           & Preemptive \checkmark \\
    Random      & $> 1$                   & $> 1$                    & $\leq 1$           & Preemptive \checkmark \\
    \bottomrule
  \end{tabular}
\end{table}

ワークロード分布によらず，プリエンプティブ方式のみが一貫して高性能ユーザを保護し，参加インセンティブを維持できることが示された．

% ============================================================
\section{考察}
% ============================================================

\subsection{プリエンプティブ方式が機能する理由}

本評価を通じて，プリエンプティブ方式のみが高性能GPU所有者の参加を維持し，システムを持続可能にできることが示された．
その本質的な理由は以下のように整理できる．

FCFS・所有者優先方式では，ゲストタスクが実行中である間，所有者がタスクを投入しても即座に処理を開始できない．
高性能GPUに対するゲスト需要は高いため，高負荷時にはGPUが常にゲストタスクで埋まった状態が続く．
結果として，所有者のTATは非共有時（自GPU単独利用）より悪化し，PRが1を超える．
合理的なユーザであれば，この状況では「参加しない方が得」と判断し，ACPから離脱する．

一方，プリエンプティブ方式では，所有者がタスクを投入した瞬間にゲストタスクが即座に停止される．
所有者の待ち時間は，自身の先行タスク（既に投入されていたタスク）が完了するまでの時間のみであり，
これは非共有環境と本質的に同一である．さらに，他者の遊休GPUをゲストとして利用できるため，
自身のキューが短い時は他機に投入することでTATを短縮できる．
すなわち，プリエンプティブ方式では「参加しても損しない，むしろ得をする」状態が保証される．

\subsection{高負荷時のプリエンプティブ方式における全体TATの上昇}

第4.2節で示した通り，$\rho_\mathrm{sys} = 0.9$以上の極高負荷域では，
プリエンプティブ方式の\textbf{全体平均}TATが所有者優先方式を上回る場合がある．
しかしこれは，低・中性能ユーザのゲスト利用時における頻繁なプリエンプションによるペナルティ蓄積が原因であり，
高性能ユーザ自身のTATは全負荷帯を通じて非共有時と同等以下に保たれている（第4.3節）．

この点は，ACPが「参加インセンティブの維持」を第一目的とするシステムであることを踏まえると本質的な問題ではない．
スループット最大化を目的とするシステムであれば，高負荷時の全体TATが上昇することは問題となるが，
ACPが目指すのは「参加者が合理的に判断した際に必ず参加を選択するような環境の構築」であり，
その目標はプリエンプティブ方式により達成されている．

\subsection{経済的報酬なしでの持続可能性}

本研究の最も重要な知見の一つは，\textbf{複雑な経済的報酬機構を用いずとも，適切な技術的制約の導入のみでシステムの持続可能性を確保できる}という点である．

P2Pシステムにおける参加維持は，従来は評判システムや金銭的インセンティブの設計によって実現されてきた\cite{feldman2004robust}．
しかし，これらの仕組みはシステムの複雑性を増し，参加者の理解・受容のコストが高い．
プリエンプションという技術的保証は，ユーザが直感的に理解しやすく（「自分のGPUは必要な時に必ず使える」），
かつ複雑なプロトコルを必要としない点で優れたインセンティブメカニズムである．

\subsection{実装上の考慮事項}

本研究はシミュレーションによる評価であり，実装に向けては以下の課題が存在する．

\textbf{プリエンプションコスト}：本研究ではオーバーヘッド係数を$\alpha = 0.2$（固定値）として扱ったが，
実際のコストはタスクの種類，チェックポイントのサイズ，ネットワーク帯域幅等に依存する．
深層学習の訓練タスクでは，GPUメモリの保存・転送が主なオーバーヘッドとなる．

\textbf{チェックポイント機構}：PyTorchやTensorFlowはモデルのチェックポイント保存機能を標準でサポートしており，
これを活用することでプリエンプション時の進捗損失を最小化できる．
ただし，完了直前のタスクに対する猶予時間（grace period）の設定等，制御の高度化が必要となる．

\textbf{マイグレーション通信オーバーヘッド}：学内ネットワークの帯域幅がマイグレーション速度を制約するため，
高速なネットワーク（10GbE以上）の整備が効果的な実装の前提条件となる．

\subsection{本研究の限界}

本研究にはいくつかの限界が存在する．
第一に，シミュレーション規模が18ユーザ・9 GPUティアという比較的小規模な環境に限定されており，
数百〜数千ユーザの大規模クラスタへのスケーラビリティは未検証である．
第二に，各ユーザが1台のGPUを所有するという単純な所有関係を前提としており，
1研究室が複数台のGPUを共同管理する構成は考慮していない．
第三に，プリエンプションコストを固定値として扱っており，タスク種別や実行段階に依存する変動コストは未モデル化である．
これらの課題については，第6章で今後の研究方向として議論する．

% ============================================================
\section{まとめと今後の課題}
% ============================================================

\subsection{本研究のまとめ}

本研究では，所有権が分散した学内環境において，資源提供者の参加インセンティブを保護しながらシステムの持続可能性を実現する，
\textbf{所有者優先権を保証するプリエンプティブ方式}を採用した学内GPU共有基盤「Air Computing Pool（ACP）」を提案・評価した．

シミュレーション評価の結果，以下の主要な成果を達成した．

\begin{itemize}
  \item \textbf{低性能層のタスク滞留解消}：
    非共有環境で数十万〜数百万秒の飽和状態にあった低性能GPU所有者のTATが，
    ACP参加により数百秒程度へと劇的に改善された．

  \item \textbf{高性能層のTAT短縮}：
    プリエンプティブ方式を用いることで，高性能GPU所有者のTATを非共有時と同等以下に維持しつつ，
    最大約78\%のTAT短縮を達成した（Tier 8, 9ユーザー対FCFS比）．

  \item \textbf{全ユーザの参加継続率100\%}：
    参加行動モデルを用いた評価により，プリエンプティブ方式では全ユーザが参加継続の判断を下し，
    第2イテレーション以降は100\%の参加率を安定して維持することが示された．
    FCFS・所有者優先方式ではいずれも高性能層の離脱によるシステム崩壊が確認されたのと対照的である．

  \item \textbf{ワークロード不均一性への頑健性}：
    均一・低性能層重視・高性能層重視・ランダムの4種類のワークロードシナリオすべてにおいて，
    プリエンプティブ方式のみが一貫してPR $\leq 1$を維持することを確認した．
\end{itemize}

以上の結果から，複雑な経済的報酬を用いずとも，適切な技術的制約（プリエンプション）の導入のみで
持続可能な共有基盤が構築可能であることを示した．

\subsection{今後の課題}

本研究の知見を実用化・発展させるための今後の課題を以下に示す．

\paragraph{実装と実環境評価}
本研究はシミュレーションによる評価であり，KubernetesやSLURMなどの実際のスケジューラへの統合実装および実環境での運用実験が今後の重要な課題である．
特に，実際のプリエンプションコストの測定と，それに基づくモデルの精緻化が必要である．

\paragraph{スケーラビリティの検証}
18ユーザ・9 GPUティアという評価環境を，数百〜数千ユーザの大規模環境へ拡張した際のシステム挙動の検証が必要である．
大規模環境では分散スケジューリングアルゴリズムの設計，スケジューラ自体のボトルネック問題，
ネットワーク遅延を考慮したモデルの拡張等が課題となる．

\paragraph{複合的なインセンティブ設計}
本研究では所有権の物理的保証のみをインセンティブとして扱ったが，
高性能GPUへの負荷集中に伴う電力消費やハードウェア消耗の問題も考慮し，
優先利用枠やポイント制，クレジット制度などを統合した多角的なインセンティブ設計の発展も必要である．

\paragraph{制御の高度化}
プリエンプションによる即時切断に対して，完了間近なタスクへの猶予時間提供や，
マイグレーション時の通信オーバーヘッド低減など，制御の高度化による実行効率の向上を検討する．

\paragraph{多様なワークロードへの対応}
現在はポアソン到着・対数正規サイズのタスクモデルを採用しているが，
現実のバースト的なタスク発生や，長時間にわたる大規模訓練ジョブ（数十〜数百時間規模）への対応についても検証が求められる．

% ============================================================
% 謝辞
% ============================================================
\section*{謝辞}

本研究はJSPS科研費JP24K14914の助成を受けたものです．
また，研究の進行にあたり，大下裕一先生，下西英之先生をはじめとする大阪大学の皆様に多大なるご指導・ご支援をいただきました．ここに深く感謝申し上げます．

% ============================================================
% 参考文献
% ============================================================
\begin{thebibliography}{99}

\bibitem{lecun2015deep}
  Y. LeCun, Y. Bengio, and G. Hinton,
  ``Deep learning,''
  \textit{Nature}, vol.521, no.7553, pp.436--444, 2015.

\bibitem{brown2020language}
  T. Brown \textit{et al.},
  ``Language Models are Few-Shot Learners,''
  \textit{Advances in Neural Information Processing Systems (NeurIPS)}, 2020.

\bibitem{hooker2021hardware}
  S. Hooker,
  ``The hardware lottery,''
  \textit{Communications of the ACM}, vol.64, no.12, pp.58--65, 2021.

\bibitem{armbrust2010view}
  M. Armbrust \textit{et al.},
  ``A view of cloud computing,''
  \textit{Communications of the ACM}, vol.53, no.4, pp.50--58, 2010.

\bibitem{buyya2009cloud}
  R. Buyya \textit{et al.},
  ``Cloud computing and emerging IT platforms: Vision, hype, and reality for delivering computing as the 5th utility,''
  \textit{Future Generation Computer Systems}, vol.25, no.6, pp.599--616, 2009.

\bibitem{jeon2019analysis}
  M. Jeon \textit{et al.},
  ``Analysis of Large-scale Multi-tenant GPU Clusters for DNN Training,''
  \textit{2019 USENIX Annual Technical Conference (ATC'19)}, 2019.

\bibitem{ghodsi2011dominant}
  A. Ghodsi \textit{et al.},
  ``Dominant Resource Fairness: Fair Allocation of Multiple Resource Types,''
  \textit{8th USENIX Symposium on Networked Systems Design and Implementation (NSDI'11)}, 2011.

\bibitem{feldman2004robust}
  M. Feldman \textit{et al.},
  ``Robust incentive techniques for peer-to-peer networks,''
  \textit{Proceedings of the 5th ACM Conference on Electronic Commerce (EC'04)}, 2004.

\end{thebibliography}

\end{document}
